\section{Referenciais}

Durante a simulação da missão de \gls{rvdb}, diferentes referenciais são utilizados para descrever as fases do movimento das espaçonaves. 
Nas etapas iniciais da missão, em que o movimento do perseguidor é descrito em relação à Terra, utiliza-se o referencial \gls{eci}. 
Nesse referencial, a dinâmica orbital é tratada como movimento absoluto da espaçonave em torno do corpo central.

Por outro lado, nas fases de aproximação de curta distância e final entre as espaçonaves, o movimento do perseguidor passa a ser descrito em relação ao alvo. 
Nesse caso, utiliza-se o referencial \gls{lvlh}, centrado no alvo, que permite representar de forma direta o movimento relativo entre as duas espaçonaves.

A Tabela~\ref{tab_fases_missao_referenciais} apresenta as fases da missão consideradas na simulação e os referenciais adotados em cada etapa.

\begin{table}[H]
\centering
\caption{Fases da missão de \gls{rvdb} e referenciais utilizados na modelagem do movimento.}
\label{tab_fases_missao_referenciais}
\begin{tabular}{lll}
\hline
\textbf{Fase da missão} & \textbf{Tipo de movimento} & \textbf{Referencial} \\
\hline
Injeção orbital                & Movimento absoluto & \gls{eci} \\
Órbita inicial                 & Movimento absoluto & \gls{eci} \\
Transferência de órbita        & Movimento absoluto & \gls{eci} \\
Aproximação de curta distância & Movimento relativo & LVLH \\
Aproximação final              & Movimento relativo & LVLH \\
\hline
\end{tabular}
\end{table}